\chapter{Belief-Flow Accounting}
\label{ch:belief-flow-accounting}

\epigraph{Follow the money.}{Deep Throat (attr.)}

\noindent
If doxonomics is to be a science rather than a commentary, it needs accounting.
The claim that ``beliefs are funded'' is empty without a ledger: who spent how much, through which institutions, producing which narratives, reaching which audiences, with what measurable effects.
This chapter develops the framework for tracing financial flows into narrative outputs: the doxonomic ledger.

Belief-flow accounting is to doxonomics what national income accounting is to macroeconomics: a systematic method for measuring aggregate flows.
Just as GDP accounting tracks the flow of goods and services through an economy, doxonomic accounting tracks the flow of resources through the belief production chain.
The analogy is not perfect (narrative production is harder to measure than automobile production), but the aspiration is the same: replace impressionistic claims with documented flows.


\section{The Accounting Framework}
\label{sec:ch14-framework}

\begin{definition}[Doxonomic Account]
\label{def:doxonomic-account}
\index{belief-flow accounting}
A \emph{doxonomic account} for belief $p$ in period $[0, T]$ is a table
\[
  \mathcal{L}_p = \{(s_k, \iota_k, c_k, \funding{k}, \Theta_k)\}_{k=1}^{m}
\]
where $s_k$ is the funding source, $\iota_k$ is the processing institution, $c_k$ is the channel, $\funding{k}$ is the expenditure, and $\Theta_k$ is the estimated doxonomic throughput.
The \emph{total doxonomic investment} in $p$ is $F_p = \sum_k \funding{k}$.
\end{definition}

A complete doxonomic account answers five questions:
\begin{enumerate}[nosep]
  \item \textbf{Who paid?} (source identification)
  \item \textbf{Who processed?} (institutional intermediation)
  \item \textbf{What was produced?} (narrative output)
  \item \textbf{Who was reached?} (audience and exposure)
  \item \textbf{What was the effect?} (belief change, behavioral consequence)
\end{enumerate}

In practice, complete accounts are rarely achievable.
Funding opacity, institutional complexity, and the difficulty of measuring belief effects mean that most doxonomic accounts are partial. Partial accounting is vastly better than no accounting.

\begin{example}[Partial Account: Climate Doubt]
\label{ex:climate-account}
\index{belief-flow accounting!climate example}
A partial doxonomic account for the belief ``climate science is uncertain'' in the U.S., 1990--2010 \autocite{oreskes2010}:

\begin{center}
\begin{tabular}{p{2.3cm}p{2.5cm}p{2.5cm}rp{1.5cm}}
\toprule
\textbf{Source} & \textbf{Institution} & \textbf{Channel} & \textbf{\$ (est.)} & \textbf{Reach} \\
\midrule
ExxonMobil & Heartland Inst. & Reports, op-eds & \$8M & Policy \\
Koch Industries & Americans for Prosperity & TV ads, rallies & \$50M+ & Mass \\
Coal industry & Global Climate Coalition & Lobbying, PR & \$13M & Congress \\
Multiple & George C. Marshall Inst. & Congressional testimony & \$2M & Elite \\
Multiple & Competitive Enterprise Inst. & Media appearances & \$5M & Broadcast \\
\bottomrule
\end{tabular}
\end{center}

This is a fragment of the total.
\textcite{oreskes2010} and subsequent investigations (Frumhoff et al., 2015; Supran and Oreskes, 2017) have documented hundreds of millions in total spending across dozens of institutions.
The account is partial because:
\begin{itemize}[nosep]
  \item not all funding is traceable (dark money, donor-advised funds),
  \item institutional overhead is not separated from narrative production,
  \item throughput estimates are rough (how many people were actually reached and influenced?),
  \item the causal link from spending to belief change is not established by the account alone (that requires Chapter~\ref{ch:causal-inference}).
\end{itemize}
\end{example}


\section{Source Identification}
\label{sec:ch14-sources}

Tracing funding to belief requires navigating layers of intermediation.
A typical chain:
\[
  \text{Corporation} \to \text{Foundation} \to \text{Donor-advised fund} \to \text{Think tank} \to \text{Publication}
\]

Each intermediary adds opacity.
The corporation may be legally required to disclose its contributions to the foundation, but the foundation may route grants through a donor-advised fund (which does not disclose the ultimate recipient), and the think tank may not disclose its funders.

\begin{definition}[Funding Chain]
\label{def:funding-chain}
\index{funding chain}
A \emph{funding chain} for belief $p$ is a directed path in the financial network:
\[
  s_0 \xrightarrow{f_1} s_1 \xrightarrow{f_2} \cdots \xrightarrow{f_l} \iota \xrightarrow{\text{produces}} n_p
\]
where $s_0$ is the ultimate source, $s_1, \ldots, s_{l-1}$ are intermediaries, $\iota$ is the narrative-producing institution, and $f_i$ is the flow at each stage.
The \emph{chain length} $l$ measures the number of intermediation layers.
Longer chains are more opaque.
\end{definition}

\subsection{Data Sources for Source Identification}

\begin{itemize}[nosep]
  \item \textbf{IRS Form 990 filings} (U.S.): tax-exempt organizations must file annual returns listing revenues, expenditures, and grants made.
  Available through ProPublica Nonprofit Explorer and Guidestar.
  \textit{Limitation}: does not require disclosure of donors to 501(c)(4) organizations; donor-advised funds disclose the fund sponsor but not the donor.

  \item \textbf{Lobbying disclosure records}: the U.S.\ Lobbying Disclosure Act requires registration of lobbyists and quarterly reports of expenditures.
  Available through OpenSecrets.
  \textit{Limitation}: covers only direct lobbying, not grassroots campaigns, astroturfing, or think-tank-mediated advocacy.

  \item \textbf{Corporate annual reports and SEC filings}: publicly traded companies disclose political spending and trade-association memberships.
  \textit{Limitation}: disclosure requirements vary; many expenditures are classified as ``education,'' ``research,'' or ``community engagement'' without specifying ideological content.

  \item \textbf{Investigative journalism databases}: ProPublica, InfluenceWatch, LittleSis, SourceWatch.
  \textit{Limitation}: coverage is selective (focused on controversial or politically salient cases); data quality varies.

  \item \textbf{Academic forensic accounting}: scholars have reconstructed funding networks through painstaking archival work.
  \textcite{mayer2016} traced Koch network funding; \textcite{oreskes2010} traced fossil fuel funding of doubt.
  \textit{Limitation}: labor-intensive, not systematic, not reproducible at scale.
\end{itemize}

\begin{definition}[Funding Opacity]
\label{def:opacity}
\index{funding opacity}
The \emph{opacity} of a funding chain is a measure of how difficult it is to trace the chain from narrative output back to ultimate source.
We define:
\[
  \text{Opacity}(\text{chain}) = 1 - \prod_{i=1}^{l} d_i
\]
where $d_i \in [0,1]$ is the disclosure probability at intermediary $i$ (the probability that the flow $f_i$ is publicly documented).
A fully transparent chain (all $d_i = 1$) has opacity 0.
A chain with one perfectly opaque intermediary ($d_i = 0$ for some $i$) has opacity 1.
\end{definition}

\begin{remark}[The Opacity Arms Race]
\label{rem:opacity-arms-race}
As doxonomic accounting improves, funders develop more sophisticated opacity mechanisms: pass-through entities, foreign intermediaries, anonymized cryptocurrency transfers, and corporate structures designed to frustrate tracing.
Doxonomic accounting is therefore an adversarial exercise: the accounting methodology must evolve in response to evasion strategies, much as tax enforcement evolves in response to tax avoidance.
\end{remark}


\section{Ideological Input-Output Tables}
\label{sec:ch14-leontief}

The belief production chain involves institutions that are both producers and consumers of narrative.
A think tank produces reports that are consumed by media; media produces coverage that is consumed by politicians; politicians produce policies that are consumed by think tanks as evidence for their next round of reports.
This circular flow is captured by the input-output framework.

\begin{model}[Doxonomic Input-Output Table]
\label{mod:io-table}
\index{input-output table}
Let $n$ institutional sectors produce narrative output.
The input-output matrix $\bm{A}$ has entries $a_{ij} \in [0,1]$ representing the fraction of institution $j$'s narrative output consumed as input by institution $i$.
The total output vector $\bm{n}$ satisfies
\[
  \bm{n} = \bm{A}\bm{n} + \bm{d}
\]
where $\bm{d}$ is the exogenous injection (direct funder input).
The solution, when $(\bm{I} - \bm{A})$ is invertible and the spectral radius of $\bm{A}$ is less than 1, is
\[
  \bm{n} = (\bm{I} - \bm{A})^{-1}\bm{d}
\]
The matrix $(\bm{I} - \bm{A})^{-1}$ is the \emph{doxonomic multiplier}\index{doxonomic multiplier}: it shows how a unit injection of narrative through one sector amplifies through the system.
\end{model}

\begin{example}[Constructing an Input-Output Table]
\label{ex:io-construction}
Consider four sectors: think tanks ($T$), prestige media ($M$), broadcast news ($B$), and education ($E$).
Estimating the inter-sector flows from citation patterns, sourcing data, and content analysis:

\begin{center}
\begin{tabular}{l|cccc|c}
\toprule
\textbf{To $\downarrow$ / From $\to$} & $T$ & $M$ & $B$ & $E$ & Exog.\ $\bm{d}$ \\
\midrule
Think tanks $T$ & 0 & 0.05 & 0.02 & 0.10 & \$50M \\
Prestige media $M$ & 0.30 & 0 & 0.10 & 0.05 & \$20M \\
Broadcast news $B$ & 0.15 & 0.40 & 0 & 0.02 & \$10M \\
Education $E$ & 0.20 & 0.10 & 0.05 & 0 & \$30M \\
\bottomrule
\end{tabular}
\end{center}

Reading the table: 30\% of think-tank output is picked up by prestige media ($a_{MT} = 0.30$); 40\% of prestige media output is recycled into broadcast news ($a_{BM} = 0.40$); 20\% of think-tank output eventually enters educational materials ($a_{ET} = 0.20$).

The Leontief inverse $(\bm{I} - \bm{A})^{-1}$ reveals:
\begin{itemize}[nosep]
  \item a \$1M injection into think tanks produces approximately \$1.5M in total narrative output (50\% amplification through the system),
  \item the largest multiplier path is $T \to M \to B$: think-tank research is picked up by prestige media, which is then amplified by broadcast news,
  \item the education sector is a ``slow amplifier'': it absorbs think-tank and media output and releases it over decades through curricula.
\end{itemize}
\end{example}

\begin{remark}[Estimating the $\bm{A}$ Matrix]
\label{rem:estimating-A}
The inter-institutional flow coefficients $a_{ij}$ are not directly observable.
They must be estimated from:
\begin{itemize}[nosep]
  \item \textbf{Citation analysis}: how often does institution $i$ cite or reference institution $j$'s outputs?
  \item \textbf{Sourcing studies}: content analyses of media output identifying the institutional origins of claims, data, and expert quotes.
  \item \textbf{Network analysis}: co-membership, co-citation, and personnel-flow networks that trace idea movement between institutions (Chapter~\ref{ch:network-analysis}).
  \item \textbf{Survey of practitioners}: journalists report which sources they rely on; policymakers report which think-tank outputs influence their thinking.
\end{itemize}
All methods are imprecise.
The input-output table should be accompanied by uncertainty bounds on each coefficient.
\end{remark}


\section{Accounting Identities}
\label{sec:ch14-identities}

\begin{principle}[Doxonomic Accounting Identity]
\label{prin:accounting-identity}
\index{accounting identity}
For each institution $\iota$ in steady state:
\[
  \text{Total funding in} = \text{Narrative output} \times \text{Unit cost} + \text{Institutional overhead}
\]
In symbols: $\funding{\iota} = \Theta_{\iota} \cdot c_{\text{unit}} + O_{\iota}$.
The \emph{doxonomic efficiency} is $\eta_{\iota} = \Theta_{\iota} / \funding{\iota} = 1 - O_{\iota}/\funding{\iota}$.
\end{principle}

Doxonomic efficiency varies widely across institution types:
\begin{itemize}[nosep]
  \item \textbf{High efficiency} ($\eta > 0.6$): small advocacy organizations, online media outlets, social media campaigns.
  Low overhead, direct narrative output.
  \item \textbf{Medium efficiency} ($\eta \approx 0.3$--$0.5$): think tanks, PR firms.
  Moderate overhead (buildings, staff, credibility maintenance), substantial narrative output.
  \item \textbf{Low efficiency} ($\eta < 0.2$): universities, government agencies.
  High overhead (research infrastructure, administration, teaching, non-doxonomic functions), but the doxonomic output has very high credibility and persistence.
\end{itemize}

\begin{remark}[Efficiency vs.\ Effectiveness]
\label{rem:efficiency-effectiveness}
Low doxonomic efficiency does not mean low doxonomic effectiveness.
A university with $\eta = 0.1$ spends 90\% of its budget on non-doxonomic functions (research, teaching, administration), but the 10\% that produces belief effects has $\credibility{} \approx 0.9$ and $\persistence{} \approx 30$ years.
A social media campaign with $\eta = 0.8$ converts most of its budget to reach, but with $\credibility{} = 0.15$ and $\persistence{} = 0.01$ years.
The long-run doxonomic impact (throughput $\times$ persistence $\times$ credibility) may be higher for the university despite its lower efficiency.
This is why the equimarginal principle (Proposition~\ref{prop:optimal-allocation}) and the IROI calculation (Definition~\ref{def:iroi}) account for credibility and persistence, not just efficiency.
\end{remark}


\section{Uncertainty in Doxonomic Accounting}
\label{sec:ch14-uncertainty}

Doxonomic accounts are never precise.
Every entry involves estimation, and honest accounting must report the uncertainty.

\begin{definition}[Account Uncertainty]
\label{def:account-uncertainty}
\index{accounting uncertainty}
For each line in a doxonomic account, the uncertainty is characterized by:
\begin{itemize}[nosep]
  \item \textbf{Source uncertainty}: the ultimate funder may be unknown or partially traced (opacity $> 0$).
  Report as a confidence interval: $s_k \in \{$known sources$\}$ with probability $p_s$, or ``unknown'' with probability $1-p_s$.
  \item \textbf{Amount uncertainty}: the expenditure $\funding{k}$ is often estimated from partial data (e.g., grant totals that include both doxonomic and non-doxonomic spending).
  Report as $\funding{k} \in [F_L, F_U]$ with 90\% confidence.
  \item \textbf{Throughput uncertainty}: the estimated doxonomic throughput $\Theta_k$ depends on reach, credibility, and persistence estimates, each with its own error.
  Report the propagated uncertainty.
  \item \textbf{Attribution uncertainty}: the fraction of an institution's output attributable to belief $p$ (vs.\ other beliefs or non-doxonomic functions) is a judgment call.
  Report the sensitivity of $F_p$ to this attribution fraction.
\end{itemize}
\end{definition}

\begin{principle}[Honest Accounting]
\label{prin:honest-accounting}
A doxonomic account without uncertainty bounds is not a doxonomic account.
It is advocacy with numbers.
The field's credibility (its own epistemic capital) depends on honest reporting of what is known, what is estimated, and what is guessed.
\end{principle}


\section{A Worked Example: Accounting for Austerity Belief}
\label{sec:ch14-austerity}

To illustrate the full accounting framework, consider a partial account for the belief ``government debt is dangerous and austerity is necessary'' in the UK, 2009--2013.

\textbf{Step 1: Source identification.}
The austerity narrative was promoted by:
\begin{itemize}[nosep]
  \item Centre for Policy Studies (CPS): founded by Thatcher and Keith Joseph, annual budget $\sim$\pounds2M.
  \item Institute of Economic Affairs (IEA): annual budget $\sim$\pounds3M.
  \item Taxpayers' Alliance: annual budget $\sim$\pounds1M.
  \item The Treasury itself: the largest single source of austerity framing, but with public rather than private funding.
  \item Financial press (\textit{Financial Times}, \textit{Economist}): editorial stances aligned with austerity, funded by subscriptions and advertising from the financial sector.
\end{itemize}

\textbf{Step 2: Processing.}
Think tanks produced policy briefs, opinion pieces, and parliamentary testimony.
The Treasury produced budget documents, forecasts, and press conferences.
Financial media produced daily coverage framing debt as crisis.

\textbf{Step 3: Throughput estimation.}
\begin{center}
\begin{tabular}{lrrcc}
\toprule
\textbf{Source} & \textbf{\pounds/yr} & \textbf{Reach} & \textbf{Credibility} & \textbf{$\Theta$} \\
\midrule
CPS & \pounds2M & 0.01 & 0.6 & 12K \\
IEA & \pounds3M & 0.01 & 0.7 & 21K \\
Taxpayers' Alliance & \pounds1M & 0.05 & 0.3 & 15K \\
Treasury (austerity share) & \pounds10M est. & 0.3 & 0.8 & 2.4M \\
Financial press & \pounds50M est. & 0.2 & 0.7 & 7M \\
\bottomrule
\end{tabular}
\end{center}

\textbf{Step 4: Total.}
The private think-tank investment was modest ($\sim$\pounds6M/year), but the amplification through the Treasury and financial press produced throughput orders of magnitude larger.
The multiplier effect (Model~\ref{mod:io-table}) was enormous: a small think-tank injection was amplified by the most credible institutions in the system.

\textbf{Step 5: Uncertainty.}
The Treasury's ``austerity share'' is an estimate (what fraction of Treasury communication was devoted to the austerity narrative vs.\ other policy); the financial press estimate is rough (advertising revenue is not doxonomic spending, but editorial alignment is).
The total investment is somewhere in the range \pounds50M--\pounds100M/year, with enormous amplification from institutional credibility.

\begin{remark}
This example illustrates a general pattern: private doxonomic investment is often small relative to the total narrative output, because public institutions and commercial media provide massive amplification.
The think tanks do not need to reach the public directly; they need to reach the \emph{institutions} that reach the public.
The multiplier, not the direct investment, is the story.
\end{remark}


\section{Accounting for Non-Monetary Flows}
\label{sec:ch14-nonmonetary}

Not all doxonomic resources are monetary.
Several important flows resist financial accounting:

\begin{itemize}[nosep]
  \item \textbf{Volunteer labor}: social movements, religious organizations, and community groups produce narrative through unpaid work.
  The monetary equivalent can be estimated (hours $\times$ wage rate), but the institutional dynamics differ from funded production.
  \item \textbf{Attention as currency}: on social media, sharing, liking, and commenting are forms of doxonomic production with zero monetary cost to the producer but significant narrative effect.
  \item \textbf{Institutional cross-subsidy}: a corporation's marketing budget (spent on product advertising) also reinforces the consumer-identity anthropological regime; a military's recruitment budget also reinforces nationalist narratives.
  These doxonomic effects are by-products, not intentional investments, but they are real.
  \item \textbf{Architecture and space}: the built environment communicates beliefs without any ongoing narrative production (monumental state buildings communicate power; open-plan offices communicate transparency/surveillance; shopping malls communicate consumption-as-activity).
\end{itemize}

\begin{definition}[Expanded Doxonomic Account]
\label{def:expanded-account}
\index{belief-flow accounting!expanded}
An \emph{expanded doxonomic account} adds non-monetary flows to the standard account:
\[
  \mathcal{L}_p^+ = \mathcal{L}_p \cup \{(\text{source}, \text{type}, \text{estimated equivalent}, \Theta)\}_{\text{non-monetary}}
\]
The ``estimated equivalent'' converts non-monetary contributions to dollar values where possible and flags them as imputed rather than observed.
\end{definition}


\section{Notes and References}
\label{sec:ch14-notes}

The accounting approach draws on national income accounting methods \autocite{stiglitz2012}.
The Leontief input-output framework originates in Leontief (1941); for a modern introduction, see Miller and Blair (2009).
Dark money flows are documented in \textcite{mayer2016}.
The application of input-output accounting to narrative flows is original to this text.
For lobbying data, see OpenSecrets; for think-tank funding, see InfluenceWatch and ProPublica Nonprofit Explorer.
On media ownership structures, see \textcite{bagdikian2004}.

On the specific case of climate-doubt funding, see \textcite{oreskes2010} and Supran and Oreskes (2017).
For the UK austerity case, see \textcite{streeck2016} and Blyth (2013), \textit{Austerity: The History of a Dangerous Idea}.
On non-monetary doxonomic production (architecture, space, material culture), see Latour (2005) and \textcite{scott1998}.
The opacity concept connects to the literature on dark money and political finance regulation \autocite{mayer2016}.


\section{Exercises}

\begin{exercise}
Using publicly available IRS 990 data (via ProPublica Nonprofit Explorer), construct a partial doxonomic account for the belief ``regulation harms economic growth'' by tracing funding from three major foundations to think tanks to publications.
Report the funding chain length and opacity for each path.
\end{exercise}

\begin{exercise}
Compute the Leontief inverse for the four-sector system in Example~\ref{ex:io-construction}.
What is the total system-wide narrative output from a \$1M injection into think tanks?
Which indirect pathway produces the largest multiplier effect?
\end{exercise}

\begin{exercise}
Estimate the doxonomic efficiency $\eta$ of: (a)~a university economics department, (b)~a social media campaign, (c)~a think tank.
For each, estimate the total budget, the fraction devoted to doxonomic output, and the resulting $\eta$.
Which has the highest efficiency?
Which has the highest long-run effectiveness?
Why might these differ?
\end{exercise}

\begin{exercise}
\label{ex:ch14-uncertainty}
Take the austerity worked example (Section~\ref{sec:ch14-austerity}) and construct 90\% confidence intervals for each throughput estimate.
How does the uncertainty in the ``austerity share'' of Treasury communication affect the total investment estimate?
Perform a sensitivity analysis: report $F_p$ under optimistic, central, and pessimistic attribution assumptions.
\end{exercise}

\begin{exercise}
\label{ex:ch14-nonmonetary}
Identify three non-monetary doxonomic flows in your daily life (e.g., architectural, spatial, routine, performative).
Estimate their doxonomic throughput in dollar-equivalent terms.
How do they compare in magnitude to the monetary flows you encounter through media and advertising?
\end{exercise}

\begin{exercise}
Construct a doxonomic account for a belief of your choice.
Report: (a)~at least three funding sources with amounts, (b)~at least two processing institutions, (c)~at least two channels with reach and credibility estimates, (d)~total $F_p$ with uncertainty bounds, (e)~a candid assessment of what you could not trace and why.
\end{exercise}

\begin{exercise}
\label{ex:ch14-opacity}
The chapter defines funding opacity as $1 - \prod_i d_i$.
For a chain of length $l = 4$ where each intermediary has disclosure probability $d_i = 0.7$, compute the opacity.
What opacity would the chain have if a donor-advised fund ($d = 0.1$) were inserted as one intermediary?
What does this imply about the doxonomic accounting challenge posed by anonymous philanthropy?
\end{exercise}
